\chapter{Filesystem}

This chapter discusses the filesystem project.
We will first analyze the performance of the provided block driver and consider
a simple optimization.
Then we lay out the design of the Filesystem infrastructure that we built on top
of this driver discussing each component in detail.

\section{Block Driver Performance}

As the book states, the provided block driver is incredibly slow.
As seen in \autoref{fig:fs_sdhc} the provided driver takes 953.55ms to read and
960.01ms to write a block from or to the SD card.
These measurements were done by repeatedly (30 times) transferring a block (512 Bytes) between a sector on the device and a fixed physical address on the host.
Reading and writing were measured separately.
This was tested with both the provided SD card as well as a model from SanDisk.
Both yielded comparable results, so that we only show the former since it was provided by the course.

Note, that this benchmark was performed while the full OS was running.
It was also executed on core 0, where most services -including the filesystem server- run.
These measurements thus represent a more realistic scenario than an isolated 
benchmark since it also measures interference from, for example, the scheduler.
This shows the performance impact of the various \mintinline{C}{barrelfish_usleep}
operations that the driver performs since these will result in yielding the
calling process.

When talking to other teams, we discovered that their baseline was more in the
range of 450ms. They appeared to be measuring performance in a more isolated setting.
While this is good for obtaining the best case performance of the driver, 
the results of our benchmark show the added performance impact of yielding in
the setting of a contested scheduler.

\begin{figure}[ht]
    \centering
    \scalebox{0.6}{\includesvg{./figs/fs_sdhc.svg} }
    \caption{SDHC Block driver performance with and without our optimization.}
    \label{fig:fs_sdhc}
\end{figure}

There is plenty of potential for speedup here, since modern SD card can reach transfer-speeds of multiple MB/s. 

We were restricted on time and thus chose not to explore the optimization of the
driver in great detail.
Instead we focus on the engineering challenge of creating a filesystem around the block device.
The filesystem will have to try to mitigate the performance impact of the slow driver.

However, We did find a small optimization that results in a roughly $2x$
speedup.

Looking at the implementation of \mintinline{c}{sdhc_read_block} and \mintinline{c}{sdhc_write_block} we noticed that a set block length command is send before every read and write.
This seemed redundant since the block length does not change during the
execution.
Removing these commands resulted in a read latency of 473.54ms per block and a
write latency of 480ms per block.
We still issue the command to set the block size once at the beginning when initializing the device.

The question remains whether this is a legitimate optimization or if we are sacrificing correctness and conformity to the protocol by doing this.
After extensive testing we did not encounter any errors when transferring blocks and thus assume that this optimization is legitimate. 

The benchmark for the block operations can be found in \mintinline{c}{usr/drivers/fsfat32/benchmarks/sd_card_bench.h}

In the following we will assume the block driver runs with our optimization.

\section{FAT32 Filesystem}

As suggested in the courses accompanying book we implement a filesystem for FAT32 formatted block devices.
The implementation tries to be close to the standard, but cuts some corners to save time.
The filesystem infrastructure consists of multiple layers that provide increasingly abstract interfaces for reading and writing to the FAT32 formatted block device.

These components are, in increasing abstraction:
\begin{itemize}
    \item An \textit{SDHC Wrapper} Library around the provided driver.
    \item A \textit{FAT32 Library} for basic file operations.
    \item A \textit{Filesystem Server} providing an RPC interface to the files and controlling access to the same.
    \item A \textit{Client Side Library (fatfs lib)} to interact with the RPC Server and maintain the state of files in the client.
\end{itemize}

Figure \ref{fig:fs_arch} gives an overview of the Architecure of the Filesystem
infrastructure.
In the following we will discuss these components in detail.

\begin{figure}[ht]
    \centering
    \scalebox{0.6}{\includesvg{./figs/FSServer_Infra.svg} }
    \caption{Architecture of the Filesystem Infrastructure.}
    \label{fig:fs_arch}
\end{figure}


\subsection{SDHC Wrapper}

The provided driver only exposes an interface for transferring single blocks between the device and a physical memory location.
Using this interface directly is clumsy which is why we create a wrapper around it. 
This wrapper provides functions for reading and writing multiple blocks and just parts of blocks to either physical or virtual addresses:

\begin{minted}{c}
errval_t read_from_sd_offset(struct sd_wrapper *sd, void *dst, uint32_t sector, 
                                uint32_t offset, uint32_t size);
errval_t write_to_sd_offset(struct sd_wrapper *sd, uint32_t sector, 
                                uint32_t offset, const void *src, uint32_t size);

errval_t read_from_sd_to_paddr(struct sd_wrapper *sd, lpaddr_t paddr, 
                                uint32_t sector, uint32_t size);
errval_t write_to_sd_from_paddr(struct sd_wrapper *sd, lpaddr_t paddr, 
                                uint32_t sector, uint32_t size);
\end{minted}

The first two functions allow reading and writing data of arbitrary sizes at any offset withing a block.

We will use the terms sector and block interchangeably. For both we assume 512 bytes.

When using \mintinline{C}{read_from_sd_offset()}, the block driver writes to an internal buffer. 
The desired chunks are then copied over to the buffer provided by the user.
Writing works analogously.

The latter two calls do not buffer and thus require less copying. 
This results in them being slightly more restrictive, as they do not allow an offset and the size parameter has to be a multiple of the block size.

\subsubsection{Idea for Optimization: Caching Blocks}
As described later in the chapter, we already have a write-through cache for FAT entries.
When writing back the entries, this uses the buffered operations provided by this wrapper.

By using a write-back cache, we could further cache the write accesses in the buffer and only write them to disk, when we need the buffer space for something else.
This would eliminates a lot of accesses to the device, since usually (assuming locality) fat entries are grouped together.

This could also be useful when iterating through directories.

It is probably less useful for writing and reading large chunks of data. The
cache will constantly be evicted by the next chunk.

This is currently not implemented, but might be an interesting investigation for
future work.

\subsection{FAT32 Library "usr/drivers/fsfat32/fat32\_file\_system.h"}

Our FAT32 provides the functionality to lookup, create, read, and modify both files and directories.
It exposes the following API:

\begin{minted}{c}
// Create and Delete Files or Directories.
errval_t fat32_create_file(struct fat32_file_system *fat, struct dir_entry *parent,
                            char *name, uint8_t attr, struct dir_entry *ret);
errval_t fat32_delete_file(struct fat32_file_system *fat, struct dir_entry *parent, 
                            struct dir_entry *dir_entry);

// Read / Write / Append files.
errval_t fat32_read(struct fat32_file_system *fat, struct dir_entry *dir_entry,
                            void *buf, size_t pos, size_t bytes, size_t *ret_bytes);
errval_t fat32_write(struct fat32_file_system *fat, struct dir_entry *dir_entry,
                            void *buf, size_t pos, size_t bytes, size_t *ret_bytes);

// Uses read_from_sd_to_paddr to directly read the entire file into a frame.
errval_t fat32_readfile_to_frame(struct fat32_file_system *fat, struct dir_entry *dir_entry,
                                struct capref frame, size_t *ret_bytes); 
// Truncate a file.
errval_t fat32_trunc(struct fat32_file_system *fat, struct dir_entry *dir_entry, size_t bytes);

// Find a file or directory by path.
errval_t fat32_resolve_path(struct fat32_file_system *fat, const char *path,
                            struct dir_entry *result);
// Find a file or directory in parent by name.
errval_t fat32_find_dirent(struct fat32_file_system *fat, struct dir_entry *parent,
                            char *filename, struct dir_entry *res);


// Finds the position and name of the next non-empty directory entry in parent after the given offset.
// Return FS_ERR_INDEX_BOUNDS if the end of the directory is reached.
errval_t read_next_dir_after_offset(struct fat32_file_system *fat, struct dir_entry *parent,
                                    uint64_t offset, char **ret_name, uint64_t *ret_idx);
\end{minted}

Files and directories are represented as dir\_entry structures which contain a copy of the FAT32 directory entry as well as the position of the directory entry on the disk.
In general, files and directories are treated interchangeably. 
They can be distinguished by the directory bit in the attribute entry.

Not all operations are allowed on both directory entry types and will return FS\_ERR\_NOTFILE or FS\_ERR\_NOTDIR as appropriate.
In general, lookup (find/resolve), creation, and deletion are defined for both.
Reading, writing, and truncating are only allowed on files.
To iterate through directories, we provide read\_next\_dir\_after\_offset.

Note that the create function will not check if a file or directory with that name already exists.
This must be checked by the caller beforehand.

Delete will only remove directories if they are empty.
We do not provide a direct interface to recursively delete a directory, since the function would have to check if files are still being used.
This is out of scope for this layer (In fact, this is completely left up to the end user of the library).

This layer notably provides an offset based layer.
Streaming functionality must be implemented on top.
Unfortunately, this means that, for every operation, files have to be traversed from the beginning until the specified offset is reached.
We made this choice to reduce the amount of state that has to be kept by the filesystem server.
A streaming API would require the server to keep track of the position in the file for every user that opens it.
Keeping in line with the design of Barrelfish, we instead outsource this responsibility to the user.

This design means that fast access to the file allocation table is critical.
We achieve this using a fat cache.

\subsubsection{Optimization: File Allocation Table Cache}
Originally, all operation in the library would read and write to the block device to obtain data.
Every read and write to an entry entry in the file allocation table resulted in a complete block operation to the SD card.
This meant that traversing files was slow.

To solve this we simply cache the file allocation table in memory.
This speeds up file traversal significantly. 

Figure \ref{fig:fs_read_one} shows the latency of reading a single cluster (4096Bytes) with a cached FAT entry vs a non cached FAT entry.


\begin{figure}[ht]
    \centering
    \scalebox{0.6}{\includesvg{./figs/fs_read_one_cluster.svg} }
    \caption{}
    \label{fig:fs_read_one}
\end{figure}

% uncached read ms     2144.038829
% cached read ms    1920.049408

However, since reading from the SD Card is prohibitively slow it is impractical to read the entire file allocation table upon startup.
This would mean a significant delay when starting up the system.
It is also most likely unnecessary since not all FAT entries are needed right away. 

Instead we load the FAT on demand.
Whenever a non cached entry is read, we read the entire cluster containing the FAT entry into the cache.
To keep consistency between the cache and the block device we use the write-through strategy.
This means that updating a FAT entry results in writing the surrounding block to the block device.
As a result, updating a fat entry is still expensive, but there is a noticeable difference reading is significantly faster.

\subsubsection{Performance}

For brevity, we evaluate the performance of a subset of the provided operations.
We measure each execution 30 times and report the median.
The variance is shown using small bars, but is usually so negligible that it is
not visible.

\begin{figure}[ht]
    \centering
         \subfloat[\centering ]{{\scalebox{0.4}{\includesvg{./figs/fs_fat32_lib_ops.svg} }}}
         \subfloat[\centering ]{{\scalebox{0.4}{\includesvg{./figs/fs_fat32_lib_ops_slow.svg} }}}
    \caption{}
    \label{fig:fs_fat32_lib_ops}
    
% creation_ms                      2.391705
% append_ms                        2.888078
% find_ms                          0.240011
% deletion_ms                     31.680818
% reading_large_file_ms_40961B    19.444631
\end{figure}

Figure \ref{fig:fs_fat32_lib_ops} clearly shows the limitations of the block
driver.
Finding a file in a short directory, creating a file, and appending to an empty
file are all relatively fast at $0.24ms$, $2.39ms$, and $2.89ms$ respectively.
All of these instructions need a limited amount of block operations.

Deleting and reading a large (40kB) file are considerably slower with $31.68ms$
and $19.44ms$.
These require more interaction with the block device.
Deleting a file takes such a long time because the operation tries to cleanup
the parent directory to release potentially empty clusters.
This must be done to prevent the SD card running out of capacity without
actually storing any data. 


\subsection{Filesystem Server "usr/drivers/fsfat/"}

The filesystem server exposes file operations using an RPC interface.
It exposed interface can be found through the nameservice by looking up "FS/<mountpoint>".
Since we only provide the filesystem for the SD card, we set this to be "FS/sdcard".

This allows multiple filesystems to be mounted as a sort of VFS.
This is currently, however, not supported by the client side library.

The filesystem exposes the following RPC calls:
\begin{minted}{c}
FS_OPEN                 // (path, flags) -> fd, size
FS_CREATE               // (path, flags, is_dir) -> fd
FS_CLOSE                // (fd) -> void
FS_READ                 // (fd, offset, size) -> data
FS_WRITE                // (fd, data, offset) -> size
FS_TRUNC                // (fd, size)
FS_OPENDIR              // (path) -> fd
FS_READNEXTDIR          // (fd, pos) -> filename
FS_DELETE               // (path) -> void       - Used for file & dirs
FS_READFILE_TO_FRAME    // (fd, frame_cap) -> size
\end{minted}

Like the FAT32 library that backs these RPC calls, the position in the file is
specified using offsets from the beginning of the file.

As the central point for file access to the mounted device, the server must ensure that two invariants are upheld.
Firstly, there must never be more then one simultaneous read issued to the underlying block device.
Secondly, write access to a file must always be exclusive while read accesses may overlap.
These are needed to guarantee the consistency of the filesystem.

The first invariant is trivially upheld, since the server handles all incoming requests sequentially in a single "event handler loop".

To ensure the second invariant, the server tracks which files are opened and
with which permissions.
For a client, trying to open a file, this means the following:
If the client wants to open the file with write permissions, it is only granted
access if the file is not currently opened.
If the client wants to read the file, the file must not already be opened with
write permissions.
Note, that the server does not keep track which processes have opened files or
where in the file they are. 

There can be multiple readers at a time, the server keeps track of the number of
readers accessing it, so it knows when the file is not being accessed anymore.
The following snippet shows the just described handle kept by the server:

\begin{minted}{c}
struct file_handle {
    char *path;
    struct dir_entry dir_entry;
    uint32_t num_accessors;
    bool has_writer;
};

struct server_state {
    bool up;
    struct fat32_file_system fat;
    collections_hash_table *open_files;  // A hash table of FD -> file_handle
};
\end{minted}

\subsubsection{File Descriptors}
The avid reader might have noticed that the RPC interface presented above uses
file descriptors to identify files:
Opening a file will return a file descriptor which then has to be passed back as
arguments for other operations on that file.

The server uses this file descriptor to find the file\_handle in its internal
table.
Since there should only every be at most one handle per file, the file
descriptor that is handed out is deterministically calculated from the file
name.
It is calculated as an SDBM hash of the full file path.
Collisions are resolved by comparing the file names and sequentially checking
each increasing hash value.

As a side-effect, this allows for fast lookup for file operations.

\subsubsection{Drawbacks}
However, it also has a significant drawback:
Since file descriptors are deterministic a malicious process could access a file
that was opened by another process without actually having to open it itself by
computing the file descriptor using the hash function.

This would easily circumvent the single writer policy that the server tries to
enforce.

A solution for this would be to hand out capabilities which then contain the
hash to identify the file.
Since Barrelfish uses capabilities for most things, this method would integrate 
well with the existing ecosystem.

We leave it as future work, since it was not in the scope of this project.

\subsubsection{No SD detected}
When the filesystem server is started it will try to access and setup the block
device. If this succeeds, it then registers itself as a nameservice and responds
to queries as described in this chapter.

If the setup does not succeed, because, for example, there is no SD card
present, the fileserver will still register itself with the nameservice system
and respond to queries.
However, it will respond to all queries with the error code FS\_ERR\_NOT\_FOUND.

This mechanism is a placeholder, since the client library will always try to 
connect to filesystem server upon setup.

One alternative solution would be for the client to try to connect a fixed 
number of times and then fail if it could not connect.
This is unreliable and leads to unnecessary latency when starting a new process.

A better solution would be to provide another service that maintains a list of
mounted file servers (essentially providing a VFS).
The filesystem servers would then register with this service when it successfully
set up.
And the client would query the "VFS" server to obtain the list of mounted
file systems.


\subsection{Client Side Library "fs/fatfs.h"}

The client side library forms the glue between the general libc library and the
RPC filesystem server.

Upon initialization the client library will establish a connection to the
filesystem server for "/sdcard".
The library exposes a virtually identical interface to the ramfs library on which 
it is based and provides a C wrapper for the servers RPC interface.
It exposes the following interface to libc:

\begin{minted}{c}
errval_t fatfs_open(void *st, const char *path, int flags, fatfs_handle_t *rethandle);
errval_t fatfs_create(void *st, const char *path, int flags, fatfs_handle_t *rethandle);
errval_t fatfs_remove(void *st, const char *path);
errval_t fatfs_read(void *st, fatfs_handle_t handle, void *buffer, size_t bytes,
                    size_t *bytes_read);
errval_t fatfs_read_file_to_frame(fatfs_mount_t st, const char *path, struct capref frame, size_t *bytes_read);
errval_t fatfs_write(void *st, fatfs_handle_t handle, const void *buffer, size_t bytes,
                     size_t *bytes_written);
errval_t fatfs_truncate(void *st, fatfs_handle_t handle, size_t bytes);
errval_t fatfs_tell(void *st, fatfs_handle_t handle, size_t *pos);
errval_t fatfs_stat(void *st, fatfs_handle_t inhandle, struct fs_fileinfo *info);
errval_t fatfs_seek(void *st, fatfs_handle_t handle, enum fs_seekpos whence, off_t offset);
errval_t fatfs_close(void *st, fatfs_handle_t inhandle);
errval_t fatfs_opendir(void *st, const char *path, fatfs_handle_t *rethandle);
errval_t fatfs_dir_read_next(void *st, fatfs_handle_t inhandle, char **retname,
                             struct fs_fileinfo *info);
errval_t fatfs_closedir(void *st, fatfs_handle_t dhandle);
errval_t fatfs_mkdir(void *st, const char *path);
errval_t fatfs_rmdir(void *st, const char *path);
errval_t fatfs_mount(const char *uri, fatfs_mount_t *retst);
\end{minted}

Since the server does not track the position of the cursor for each process, the
client library has to do this itself.
It keeps a list of file handles that contain basic information about the files.
This includes, the position of the cursor and the filesize.
The client library is responsible for updating its own view information when
modifying the file.
This does not introduce consistency problems, since there can only ever be one
writer per file at a time (provided there is no malicious actor as described
above).
As a result, operations like fstat, fseek, and ftell do not require an RPC.
The filesystem server only needs to be accessed to open and close files or to
access information that is on the block device.

\subsection{Performance of the complete stack}

Finally, we want to take a look at the end to end performance of issuing a
filesystem operation from userspace.

We run the benchmark from a userspace program that is spawned after the filesystem server is set up.
All measurements were run 30 times with a complete running OS.
Append and read both operate on the first cluster (4KiB) of a file.

\begin{figure}[ht]
    \centering
    \scalebox{0.6}{\includesvg{./figs/fs_fat32_e2e.svg} }
    \caption{Performance of calling libc file operations: fopen, fwrite, fread, fclose}
    \label{fig:fs_libc_e2e}
% open_ms      5494.123975
% append_ms    8656.162013
% read_ms      3840.110825
% close_ms        0.050854
\end{figure}

The results of this benchmark can be seen in figure \ref{fig:fs_libc_e2e}.
We only inspect the most prominent operations for the sake of brevity.

Opening a file takes $5494.12ms$, this is due to the fact that the call also needs to create the file.
The implementation of fopen first checks if the file exists by trying to open it.
When this fails, because the file does not exist yet, it triggers another RPC to create the file.
This explains the additional delay compared to the direct create operation.

Appending the first block to a file file takes $5494.12ms$, which includes communication overhead and traversing the directory to obtain the file handle.
It is important to note, that libc does not try to write the entire cluster in one attempt. 
Instead it breaks it into 1024 byte chunks and triggers RPCs for each.
This explains the added overhead over the direct FAT32 Library call.

Reading is subjected to the same chunking. It thus also introduces overhead resulting in a median latency of $3840.11ms$.

Closing a file multiple orders of magnitude faster at $0.051$ since it does not require any block operations.
It is just a lookup in a hashmap and an RPC.


\section{Limitations}
As can be seen in this chapter, there are still some limitation to our implementation
\begin{enumerate}
    \item The block driver is still incredibly slow
    \item Loading ELF binaries is not implemented yet, binaries are quite large and would take an unacceptably  long time to load
    \item Anything that modifies FAT entries is slow, since the block is written back every time an entry is modified.
    \item We do not support timestamps.
    \item The access control is very limited. Integrating it with the capability system would be desirable.
    \item We currently do not support long names.
    \item Currently every file open is interpreted as a write, we would just need to add flag parsing to allow simultaneous reads.
    \item Open files are not automatically closed when the process dies.
\end{enumerate}
